% ═══════════════════════════════════════════════════════════════════════
\section{L'analyse en composantes principales}
% ═══════════════════════════════════════════════════════════════════════

\begin{frame}{Le tableau centré et la matrice de covariance}
  \small $n$ individus, $p$ variables. Après avoir centré chaque variable
  (et l'avoir réduite, s'il y a lieu), on range les données dans une
  matrice $X$ de taille $n\times p$ ; la ligne $i$ est
  $\vec x_i^{\mathsf T}$. La somme des lignes est le vecteur nul.
  \begin{Enonce}{Définition 5.4 --- Matrice de covariance}
    La \defi{matrice de covariance} du tableau centré $X$ est la matrice
    $p\times p$
    \[
      C=\frac1n\,X^{\mathsf T}X .
    \]
    Son coefficient $(j,k)$ est $c_{jk}=\frac1n\sum_{i=1}^n x_{ij}x_{ik}$, la
    covariance des variables $j$ et $k$ ; ses coefficients diagonaux sont
    les variances. Si les variables ont été réduites, $C$ est la matrice de
    corrélation.
  \end{Enonce}
\end{frame}

\begin{frame}{Proposition 5.3 --- Propriétés de la matrice de covariance}
  \begin{Enonce}{Proposition 5.3 --- Propriétés de la matrice de covariance}
    Soit $C$ la matrice de covariance d'un tableau centré $X$.
    \begin{enumerate}
      \item $C$ est symétrique.
      \item Pour tout $\vec u\in\R^p$ unitaire, $\vec u^{\mathsf T}C\,\vec u$
        est la variance des projections $t_i=\vec x_i\cdot\vec u$ sur la
        droite engendrée par $\vec u$. En particulier,
        $\vec u^{\mathsf T}C\,\vec u\geq0$.
      \item Les valeurs propres de $C$ sont réelles et positives ou nulles,
        et leur somme est la \defi{variance totale}
        $c_{11}+\dots+c_{pp}$, la moyenne des carrés des distances au point
        moyen.
    \end{enumerate}
  \end{Enonce}
  \begin{AuTableau}
    \footnotesize (1) $C^{\mathsf T}=\frac1nX^{\mathsf T}X=C$. (2)
    $\vec u^{\mathsf T}C\vec u=\frac1n(X\vec u)\cdot(X\vec u)=\frac1n\sum t_i^2$.
    (3) Théorème spectral ; $\lambda=\vec v^{\mathsf T}C\vec v\geq0$ ;
    $\tr(QDQ^{\mathsf T})=\tr(D)$.
  \end{AuTableau}
\end{frame}

\begin{frame}{Remarque 5.4 --- Faut-il toujours réduire les variables ?}
  \begin{Rem}{Remarque 5.4}
    Non. Réduire, c'est donner à chaque variable la même importance au
    départ (une variance de $1$). C'est généralement souhaitable :
    \begin{itemize}
      \item lorsque les unités diffèrent (mm et g, secondes et mètres) : la
        variance totale additionnerait des quantités incomparables, et le
        résultat dépendrait des unités (exercices 5.4 et 5.9) ;
      \item lorsque les ordres de grandeur sont si différents que la plus
        dispersée imposerait le premier axe (la nageoire et le bec).
    \end{itemize}
  \end{Rem}
\end{frame}

\begin{frame}{Remarque 5.4 (suite) --- Quand ne pas réduire}
  \begin{Rem}{Remarque 5.4 (suite)}
    Même unité et dispersions comparables : les différences de variance
    font partie de ce qu'on veut décrire, et la réduction les efface.
    \begin{itemize}
      \item Des positions dans l'espace : réduire étire chaque axe, et
        déforme le nuage.
      \item Un examen où tous ont à peu près la même note : réduit, il
        pèserait autant que les autres, et ses petits écarts, dus en partie
        au hasard, seraient amplifiés.
      \item Une variable constante ne peut pas être réduite : son écart type
        est nul.
    \end{itemize}
    Les deux analyses ne donnent pas les mêmes axes : l'une diagonalise la
    matrice de covariance, l'autre la matrice de corrélation. Il faut donc
    toujours préciser laquelle on a effectuée. Dans le doute, et dès que les
    unités diffèrent, on réduit (c'est ce que font par défaut plusieurs
    logiciels).
  \end{Rem}
\end{frame}

\begin{frame}{Définition 5.5 --- Analyse en composantes principales}
  \small Théorème spectral : $C=QDQ^{\mathsf T}$, colonnes de $Q$
  orthonormées, $\lambda_1\geq\lambda_2\geq\dots\geq\lambda_p\geq0$.
  \begin{Enonce}{Définition 5.5 --- Analyse en composantes principales}
    Soit $X$ un tableau centré, $C$ sa matrice de covariance, et
    $(\vec v_1,\dots,\vec v_p)$ une base orthonormée de vecteurs propres de
    $C$, associés à $\lambda_1\geq\dots\geq\lambda_p$.
    \begin{enumerate}
      \item La droite engendrée par $\vec v_k$ est le $k$\ap{e} \defi{axe
        principal}.
      \item La $k$\ap{e} \defi{composante principale} est la variable $y_k$
        de valeurs $y_{ik}=\vec x_i\cdot\vec v_k$ ; matriciellement, $Y=XQ$.
      \item La \defi{proportion de variance expliquée} par le $k$\ap{e} axe
        est $\dfrac{\lambda_k}{\lambda_1+\dots+\lambda_p}$.
    \end{enumerate}
  \end{Enonce}
  \small Vocabulaire variable : « composantes principales » désigne parfois
  les $\vec v_k$, et les logiciels appellent les $y_{ik}$ des scores.
  L'essentiel : distinguer les directions $\vec v_k\in\R^p$ des
  coordonnées $y_{ik}$ des individus.
\end{frame}

\begin{frame}{Théorème 5.4 --- Propriétés de l'ACP}
  \begin{Enonce}{Théorème 5.4 --- Propriétés de l'ACP}
    \footnotesize
    \begin{enumerate}
      \item Chaque $y_k$ est de moyenne nulle et de variance $\lambda_k$.
      \item Deux composantes principales distinctes sont non corrélées.
      \item Pour tout $q\leq p$, parmi les sous-espaces de dimension $q$,
        $\Span\{\vec v_1,\dots,\vec v_q\}$ est celui sur lequel la variance
        du nuage projeté est la plus grande ; elle vaut
        $\lambda_1+\dots+\lambda_q$.
      \item De façon équivalente, si l'on remplace chaque individu par sa
        projection sur ce sous-espace, la moyenne des carrés des distances
        vaut $\lambda_{q+1}+\dots+\lambda_p$, et aucun autre sous-espace de
        dimension $q$ ne fait mieux.
    \end{enumerate}
  \end{Enonce}
  \begin{AuTableau}
    \scriptsize (1)--(2) $\frac1n(X\vec v_j)\cdot(X\vec v_k)=\lambda_k\,\vec v_j\cdot\vec v_k$.
    (3) Pour $q=1$ : $\vec u=\sum c_j\vec v_j$ donne
    $\vec u^{\mathsf T}C\vec u=\sum\lambda_jc_j^2\leq\lambda_1$ --- une moyenne
    des valeurs propres pondérée par des poids $c_j^2$ de somme $1$. Pour
    $q$ quelconque : des poids entre $0$ et $1$, de somme $q$. (4)
    Pythagore.
  \end{AuTableau}
\end{frame}

\begin{frame}{Des variables redondantes, des composantes qui ne le sont pas}
  \begin{itemize}
    \item Les variables de départ sont souvent fortement corrélées, donc
      redondantes : la longueur du bec et celle de la nageoire disent en
      partie la même chose.
    \item Les composantes principales sont non corrélées (point 2) : chacune
      apporte une information que les précédentes ne contenaient pas.
    \item Le point 4 mesure ce que l'on perd : la moyenne des carrés des
      erreurs de reconstruction est la somme des valeurs propres
      abandonnées.
    \item Mais c'est une moyenne : un individu particulier peut être bien
      plus mal représenté que les autres (exercice 5.3 e).
  \end{itemize}
\end{frame}

\begin{frame}{Algorithme 5.1 --- Analyse en composantes principales}
  \begin{Algo}{Algorithme 5.1 --- Analyse en composantes principales}
    Partant d'un tableau de $n$ individus et de $p$ variables :
    \begin{enumerate}
      \item centrer chaque variable ; si les unités diffèrent, ou si les
        dispersions sont très différentes, réduire aussi ;
      \item calculer la matrice de covariance $C=\frac1nX^{\mathsf T}X$ ;
      \item calculer les valeurs propres et des vecteurs propres unitaires
        de $C$, par ordre décroissant des valeurs propres ;
      \item choisir le nombre $q$ d'axes à conserver ;
      \item calculer $y_{ik}=\vec x_i\cdot\vec v_k$ et représenter le nuage
        projeté, par exemple dans le plan des deux premiers axes ;
      \item interpréter les axes à l'aide des composantes des vecteurs
        propres.
    \end{enumerate}
  \end{Algo}
\end{frame}

\begin{frame}{Combien d'axes conserver ?}
  Pas de règle universelle ; trois critères courants :
  \begin{itemize}
    \item une proportion fixée de la variance totale, par exemple
      $\SI{80}{\percent}$ ou $\SI{90}{\percent}$ ;
    \item le \defi{diagramme des éboulis} : les valeurs propres par ordre
      décroissant, et un « coude » au-delà duquel elles diminuent lentement
      (souvent du bruit) ;
    \item pour des données réduites, le \defi{critère de Kaiser} : garder les
      axes de valeur propre supérieure à $1$. La variance totale vaut $p$ ;
      un axe sous $1$ résume moins qu'une seule variable de départ.
  \end{itemize}
  \vspace{4pt}
  Pour visualiser, la question est tranchée d'avance : $q=2$ (ou $3$). Il
  faut alors garder en tête la proportion de variance conservée : elle dit
  jusqu'à quel point le dessin est fidèle.
\end{frame}

\begin{frame}{Remarque 5.5 --- Le signe des axes est arbitraire}
  \begin{Rem}{Remarque 5.5}
    \begin{itemize}
      \item Un vecteur propre unitaire peut être remplacé par son opposé
        (exercice 4.1).
      \item Changer $\vec v_k$ en $-\vec v_k$ change le signe de tous les
        $y_{ik}$ : la figure est retournée selon le $k$\ap{e} axe, sans rien
        changer à l'interprétation.
      \item Deux logiciels, ou deux versions d'un même logiciel, peuvent
        produire des figures symétriques l'une de l'autre.
      \item Dans les notes, on a choisi les signes qui rendent les figures
        plus lisibles (les papous à droite).
    \end{itemize}
  \end{Rem}
\end{frame}

\begin{frame}[fragile]{Programme 5.1 --- L'ACP avec Python}
  \begin{columns}[T]
    \begin{column}{0.66\textwidth}
\begin{lstlisting}
import numpy as np
import pandas as pd
url =("https://raw.githubusercontent.com/allisonhorst/"
       "palmerpenguins/main/inst/extdata/penguins.csv")
manchots = pd.read_csv(url).dropna(subset=["bill_length_mm"])
X = manchots[["bill_length_mm", "bill_depth_mm",
              "flipper_length_mm"]].to_numpy()
def acp(X, reduire=True):
    Z = X - X.mean(axis=0)               # centrer
    if reduire:
        Z = Z / X.std(axis=0)            # réduire : cotes z
    C = Z.T @ Z / len(Z)                 # matrice de covariance
    valeurs, vecteurs = np.linalg.eigh(C)
    ordre = np.argsort(valeurs)[::-1]    # ordre décroissant
    valeurs, vecteurs = valeurs[ordre], vecteurs[:, ordre]
    return valeurs, vecteurs, Z @ vecteurs
valeurs, vecteurs, Y = acp(X)
print(valeurs, valeurs / valeurs.sum())
\end{lstlisting}
    \end{column}
    \begin{column}{0.32\textwidth}
      \footnotesize
      \begin{itemize}
        \item \texttt{pandas} lit le fichier ; \texttt{dropna} élimine les
          deux manchots aux mesures manquantes.
        \item \texttt{np.linalg.eigh} (matrices symétriques) renvoie les
          valeurs propres par ordre croissant : on les remet dans l'ordre
          décroissant avec \texttt{np.argsort}.
        \item Sortie : valeurs propres $[\num{2.003}\ \ \num{0.767}\ \ \num{0.23}]$,
          proportions $[\num{0.668}\ \ \num{0.256}\ \ \num{0.077}]$.
        \item Le premier vecteur propre sort avec le signe opposé : les
          papous sont à gauche (remarque 5.5).
      \end{itemize}
    \end{column}
  \end{columns}
\end{frame}

\begin{frame}{Remarque 5.6 --- L'ACP en pratique}
  \begin{Rem}{Remarque 5.6}
    \begin{itemize}
      \item Les bibliothèques spécialisées, comme \texttt{scikit-learn},
        effectuent l'ACP sans programmer les calculs.
      \item Avec \texttt{PCA().fit(Z)} : l'attribut
        \texttt{explained\_variance\_ratio\_} donne les proportions de
        variance expliquée, la méthode \texttt{transform} calcule les
        composantes principales.
      \item Attention : \texttt{explained\_variance\_} divise par $n-1$
        (remarque 5.2). Les proportions, elles, ne changent pas.
    \end{itemize}
  \end{Rem}
\end{frame}
